\section{Rekursive Eigenschaft der Kugelvolumen\label{gamma:section:spheres}}
\kopfrechts{Kugelvolumen}

$\pi$ ist definiert als die Konstante, welche
\index{pi@$\pi$}%
\begin{equation}
\label{gamma-eq:circumference}
    2 \pi r = \text{Umfang eines Kreises mit dem Radius $r$}
\end{equation}
erfüllt.

Eine Kreisscheibe ist die von einem Kreisrand begrenzte, vollständig ausgefüllte Fläche eines Kreises.
\index{Kreisscheibe}%
Mit Hilfe von \eqref{gamma-eq:circumference} lässt sich die Fläche einer Kreisscheibe herleiten. 
Geht die Breite eines Rings gegen 0, verhält sich das 
Wachtstum des Flächeninhalts gleich wie das des Umfangs des Rings.
Das macht eine Herleitung der Kreisscheibenfläche per Integral möglich:
\begin{equation}
\label{gamma-eq:area}
    A(r) = \int_{0}^{r} 2 \pi r \,dr = \pi r^2
.
\end{equation}
Für eine grafische Darstellung dieser
Intuition siehe Abb. \ref{gamma-fig:ring-integration}.

Der Satz des Pythagoras gibt uns eine einfache Beschreibung aller
\index{Pythagoras}%
Punkte $(x, y)$, welche
\begin{equation}
\label{gamma-eq:pythagoras}
    x^2 + y^2 \leq r^2
\end{equation}
erfüllen und somit Teil des Kreises mit Radius $r$ sind.

Aufgrund der Ungleichung~\eqref{gamma-eq:pythagoras}
lässt sich eine weitere Definition der Kreisfläche herleiten, 
indem wir nach $y$ auflösen und danach für jedes $x$ den jeweils höchsten $y$-Wert integrieren:
\begin{gather}
\label{gamma-eq:pythagoras-y}
    y \leq \!\sqrt{r^2 - x^2} \\
\label{gamma-eq:area-pythagoras}
    A(r) = 2 \int_{-r}^{r} \!\sqrt{r^2-x^2} \,dx.
\end{gather}
Die Intuition dahinter wird in Abb.~\ref{gamma-fig:rect-integration} dargestellt.
Das Integral aus~\eqref{gamma-eq:area-pythagoras} ist nicht trivial zu lösen
und wird nicht weiter verfolgt.
Ein Vorteil ist, dass sich die durch den Satz des Pythagoras definierte Metrik
einfach auf weitere Dimensionen erweitern lässt, beispielsweise mit
\begin{equation}
\label{gamma-eq:pythagoras-3d}
    \begin{split}
    x^2 + y^2 + z^2 \leq r^2 \\
    z \leq \!\sqrt{r^2 - x^2 - y^2}
    \end{split}
\end{equation}
auf 3 Dimensionen.

Nun können wir das Volumen einer Kugel herleiten, indem wir Halbzylinder entsprechender
Grösse entlang $x$ integrieren (siehe Abb. \ref{gamma-fig:semicircle-integration}),
während iterativ in jedem Halbkreis entlang der $y$-Achse integriert wird. Das Volumen kann also mit
\begin{equation}
\label{gamma-eq:sphere-volume}
    V(r) = 2 \int_{-r}^{r} \int_{-\!\sqrt{r^2-x^2}}^{\!\sqrt{r^2-x^2}} \sqrt{r^2-x^2-y^2} \,dy \,dx
\end{equation}
berechnet werden. Eine formelle Begründung für diese iterative angehensweise ist der Satz von Fubini \cite{fubini1907integrali}.
\index{Satz!von Fubini}%
\index{Fubini, Satz von}%
Die Grenzen des inneren Integrals kommen zustande, weil für ein bestimmtes $x$ die möglichen $y$-Werte gemäss Ungleichung~\eqref{gamma-eq:pythagoras-y} eingeschränkt sind.
Für ein bestimmtes $x$ und $y$ wird dann über den maximalen Wert für $z$ gemäss Ungleichung~\eqref{gamma-eq:pythagoras-3d} integriert.

\input{papers/gamma/spheres_drawings.tex}

Da wir aber bereits die Fläche eines Kreises berechnen können, können wir $A(r)$
zur Berechnung des Kugelvolumens anstatt der Pythagoras-Methode verwenden, indem wir in
\begin{equation}
\label{gamma-eq:composed-sphere-volume}
    V(r) = 2 \int_{-r}^{r} A\left(\!\sqrt{r^2-x^2}\right) \,dx
\end{equation}
das innere Integral von~\eqref{gamma-eq:sphere-volume} durch die entsprechende Flächenberechnung ersetzen.
Das ist interessant, weil wir die Berechnung für zwei Dimensionen verwendet haben,
um das Volumen in drei Dimensionen zu berechnen.

Auf die gleiche Weise lässt sich das Volumen auf beliebige Dimensionen $> 2$ erweitern,
und zwar mit
\begin{equation}
\label{gamma-eq:n-sphere-volume-recursive}
    V_{n+1}(r) = \int_{-r}^{r} V_n\left(\!\sqrt{r^2-x^2}\right) \,dx
.
\end{equation}

\noindent
Betrachtet man $A(r) = V_2(r) = 2 \pi r$, so scheint es intuitiv,
dass $V_n$ jeweils mit $C_n r^n$ dargestellt werden kann, wobei $C_n$ eine der Dimension zugehörige Konstante ist.
Wir können induktiv beweisen, dass $V_{n+1}$ dieser Form entspricht, sofern
$V_n$ dieser Form entspricht.
Nehme an, dass $V_n(r) = C_n r^n$. Dann gilt:

\begin{align*}
%\label{gamma-eq:const-structure-induction}
         V_{n+1}(r) &= \int_{-r}^{r} C_n \left(\!\sqrt{r^2-x^2}\right)^n \,dx  \\ 
        &= \int_{-r}^{r} C_n \left(r^2-x^2\right)^{n/2} \,dx \\
        &= \int_{-1}^{1} C_n \left(r^2 - (ru)^2\right)^{n/2} r \,du \\
        &= \int_{-1}^{1} C_n \left(r^2 (1 - u^2)\right)^{n/2} r \,du \\
        &= \int_{-1}^{1} C_n r^{n+1} (1- u^2)^{n/2} \,du \\
        &= r^{n+1} \underbrace{C_n \int_{-1}^{1} (1- u^2)^{n/2} \,du}_{C_{n+1}}.
\end{align*}

\noindent
Wir sehen, dass $C_n$ ein Faktor von $C_{n+1}$ ist. Die mit dem Volumen assoziierte Konstante 
scheint also eine rekursive Struktur ähnlich der einer Fakultät aufzuweisen. Eine Untersuchung des Terms $I_n = \int_{-1}^1 (1- u^2)^{n/2} \,du$ ergibt weiter die rekursive Struktur $I_n = \frac{n}{n+1} I_{n-2}$. Dies kann durch die Rechnung
\begin{align*}
%\label{gamma-eq:n-sphere-constant}
        I_n &= \int_{-1}^1 (1- u^2)^{n/2} \,du
\\
        &= \left[ u (1-u^2)^{n/2} \right]_{-1}^{1} - \int_{-1}^1 u \cdot (-nu(1-u^2)^{n/2 - 1}) \,du
\\
        &= (0 - 0) - \int_{-1}^1 -nu^2(1-u^2)^{n/2 - 1} \,du
\\
        &= n \int_{-1}^1 u^2 (1 - u^2)^{n/2 - 1} \,du
\\
        &= n \int_{-1}^1 (1 - (1 - u^2)) (1-u^2)^{n/2 - 1} \,du
\\
        &= n \int_{-1}^1 (1 - u^2)^{n/2 - 1} \,du - n \int_{-1}^1 (1-u^2)^{n/2} \,du
\\
        &= n I_{n-2} - n I_n \\
        (n+1) I_n &= n I_{n-2} \\
        I_n &= \frac{n}{n+1} I_{n-2}
\end{align*}
gezeigt werden.

Mit diesem Wissen lässt sich der konstante Faktor des Volumens einer $n$-dimensionalen Kugel herleiten als
$C_n = C_{n-1} I_{n-1} = C_{n-2} I_{n-1} I_{n-2} = C_{n-2} \frac{2}{n} I_1 I_0$.
Der letzte Schritt ist durch
\begin{align}
    I_{n-1} I_{n-2}
&=
\biggl(\frac{n-1}{n} \cdot \frac{n-3}{n-2} \cdot \ldots \cdot \frac{3}{4}
\cdot I_1\biggr)
\cdot
\biggl(\frac{n-2}{n-1} \cdot \frac{n-4}{n-2} \cdot \ldots \cdot \frac{2}{3} \cdot I_0\biggr)
\notag\\
    &= \frac{n-1}{n} \cdot \frac{n}{n-1} \cdot \ldots \cdot \frac{3}{4} \cdot  \frac{2}{3} \cdot I_1 I_0
\label{gamma-eq:n-sphere-constant-2}
\\
    &= \frac{2}{n} \cdot I_1 I_0
\notag
\end{align}
begründet, funktioniert aber nur mit geraden $n$.

Von~\eqref{gamma-eq:area} wissen wir, dass $C_2 = \pi$.
Eine einfache Rechnung ergibt, dass $I_0 = 2$ und $I_1 = \frac{\pi}{2}$, d.h. $C_n = C_{n-2} \frac{2 \pi}{n}$.
Jetzt können wir für gerade Dimensionen die rekursive Definition `abrollen' in einen einzigen
Term:
\begin{equation}
\label{gamma-eq:nonrecursive-constant}
\begin{split}
        C_n &= \pi \cdot \frac{2 \pi}{4} \cdot \frac{2 \pi}{6} \cdot \ldots \cdot \frac{2 \pi}{n} \\
            &= \frac{2 \pi}{2} \cdot \frac{2 \pi}{4} \cdot \frac{2 \pi}{6} \cdot \ldots \cdot \frac{2 \pi}{n} \\
            &= \frac{(2 \pi)^{n/2}}{2^{n/2}} \frac{1}{\left(\frac{n}{2}\right)!} \\
            &= \frac{\pi^{n/2}}{\left(\frac{n}{2}\right)!}.
\end{split}
\end{equation}
Wir können jetzt also mit
\begin{equation}
\label{gamma-eq:n-sphere-volume-even}
    V_n(r) = \frac{\pi^{n/2}}{(\frac{n}{2})!}  r^n
\end{equation}
ein Kugelvolumen bei geraden Dimensionen einfach berechnen.

Schön wäre es, das Volumen nicht nur für gerade-dimensionale Kugeln berechnen können.
Dazu müssten wir allerdings die Fakultät auf halbe Zahlen anwenden können, was einer der Gründe ist, die Fakultät zu erweitern.
